%====================================================================
\section{\TeX\ 接口宏（\texttt{tkz-elements.sty}）}
\label{sec:sty_interface}

\subsection{用途}

文件 \texttt{tkz-elements.sty} 提供了一小批宏，作为 \TeX{} 与 \tkzNamePack{tkz-elements} 所用 Lua 引擎之间的桥梁。它们的主要目标是：

\begin{itemize}
\item 在 \TeX{} 中安全地打印 Lua 值；
\item 把几何对象（点、路径、圆）从 Lua 传递到 TikZ；
\item 在 Lua 完成计算后，简化绘图阶段。
\end{itemize}

\medskip
除非另有说明，这些宏都假定 \tkzNamePack{tkz-elements} 已在 Lua 中初始化（通常通过

|\directlua{init_elements()}|。

%--------------------------------------------------------------------
\subsection{所提供宏一览}


\begin{tabular}{ll}
\toprule
\texttt{宏} & \texttt{参考} \\
\midrule
\tkzMacro{tkz-elements}{tkzUseLua\{...\}}               & \\
\tkzMacro{tkz-elements}{tkzPrintNumber\{...\}}          & \\
\tkzMacro{tkz-elements}{tkzEraseLuaObj\{name\}}         & \\
\tkzMacro{tkz-elements}{tkzDrawCoordinates(\dots)}      & \\
\tkzMacro{tkz-elements}{tkzGetPointsFromPath\{P\}\{A\}} & \\
\tkzMacro{tkz-elements}{tkzDrawPointsFromPath(\dots)}   & \\
\tkzMacro{tkz-elements}{tkzDrawSegmentsFromPaths(\dots)}& \\
\tkzMacro{tkz-elements}{tkzDrawCirclesFromPaths(\dots)} & \\
\bottomrule
\end{tabular}

\medskip
\texttt{注意：} Lua 对象的内部布局可能会演进；这些宏在用户层面保持稳定。

%====================================================================
\subsection{宏 \tkzMacro{tkz-elements}{tkzUseLua}}
\label{subsec:tkzUseLua}

\medskip
\texttt{语法： }

\verb|tkzUseLua{<lua-expression>\}|


\medskip
\texttt{描述： } 对该 Lua 表达式求值，并把结果打印到 \TeX{} 流中。

\medskip
\texttt{参数： }
\begin{itemize}
\item \tkzname{<lua-expression>}：返回一个值的 Lua 表达式。
\end{itemize}

\medskip
\texttt{说明： }
\begin{itemize}
\item 本宏用于\emph{值}（数、布尔值、字符串）。
\item 对于 Lua 对象（点、直线、圆、路径……），请改用专用的传递/绘图宏，或调用一个显式返回值的 Lua 方法。
\end{itemize}

\texttt{示例： }
\begin{minipage}{.52\textwidth}
\begin{verbatim}
\directlua{init_elements()
  z.A = point(0,0)
  z.B = point(4,2)
}
The slope is \tkzUseLua{(z.B.im-z.A.im)/(z.B.re-z.A.re)}.
\end{verbatim}
\end{minipage}

%%%%%%%%%%%%%%%%%%%%%%%%%%%%%%%

该宏对一个 Lua 表达式求值，并把结果插入文档。


\texttt{语法： }\verb|\tkzUseLua{<lua-code>}|

\texttt{用法示例： }

\begin{mybox}
\begin{verbatim}
\tkzUseLua{utils_checknumber(math.pi)}
\end{verbatim}
\end{mybox}


%====================================================================
\subsection{宏 \tkzMacro{tkz-elements}{tkzPrintNumber}}
该宏用 PGF 的定点输出格式化并打印一个数，默认精度为 2 位小数。

\medskip

%\texttt{Syntax: }\verb|\tkzPrintNumber[<precision>]{<number>}|
\texttt{语法： }
\begin{quote}
\tkzMacro{tkz-elements}{tkzPrintNumber\{<lua-expression>\}} \qquad
\tkzMacro{tkz-elements}{tkzPN\{<lua-expression>\}}
\end{quote}

\texttt{示例： }

\begin{mybox}
\begin{verbatim}
\tkzPrintNumber{pi} % outputs 3.14
\tkzPrintNumber[4]{sqrt(2)} % outputs 1.4142
\end{verbatim}
\end{mybox}

\texttt{别名： } 宏 \tkzMacro{tkz-elements}{tkzPN} 是 \tkzMacro{tkz-elements}{tkzPrintNumber} 的简写。


\texttt{示例： }
\begin{minipage}{.52\textwidth}
\begin{verbatim}
\directlua{init_elements()
  x = math.pi/7
}
$\alpha = \tkzPN{x}$.
\end{verbatim}
\end{minipage}


%====================================================================
\subsection{宏 \tkzMacro{tkz-elements}{tkzEraseLuaObj}}

\label{sub:macro_tkzEraseLuaObj}

\texttt{语法： }
\begin{quote}
\tkzMacro{tkz-elements}{tkzEraseLuaObj\{<name>\}}
\end{quote}



\texttt{描述： }
从全局存储表中移除一个已存储的 Lua 对象（根据你的约定，通常是 \code{z}、\code{L}、\code{C}、\code{CO} 等）。


\texttt{语法： } \verb|\tkzEraseLuaObj{<object>}|

\texttt{用法示例： }

\begin{mybox}
\begin{verbatim}
\tkzEraseLuaObj{z.A}
\tkzEraseLuaObj{T.ABC}
\end{verbatim}
\end{mybox}


\texttt{说明： }
在编译大型文档、或在不同图形间复用相同对象名时，使用本宏可避免命名冲突。



%====================================================================
\subsection{宏 \tkzMacro{tkz-elements}{tkzPathCount}}
该宏取回一个 Lua 路径中的顶点个数，并把结果存入一个 TeX 宏。

\texttt{语法： }\verb|\tkzPathCount(<lua-path>){<MacroName>}|

\texttt{参数： }

\verb|<lua-path>|：一个 Lua \verb|path| 对象的标识符。
\verb|<MacroName>|：一个 TeX 宏的名字（不带反斜杠），用于存放顶点个数。


\texttt{用法示例： }

\begin{tkzexample}[latex=.5\textwidth]
\directlua{
 z.A = point(0, 0)
 z.B = point(5, 4)
 L.AB = line(z.A, z.B)
 PA.AB = L.AB:path(3)
 z.O = point(2, 0)
 }
\tkzPathCount(PA.AB){N}
\begin{center}
\begin{tikzpicture}
\tkzGetNodes
\tkzDrawCoordinates[blue](PA.AB)
\tkzDrawPointsFromPath[red,size=2](PA.AB)
  \foreach \i in {1,...,\N}{
 \tkzGetPointFromPath(PA.AB,\i){P\i}
 \tkzDrawSegment[orange](O,P\i)
  }
\end{tikzpicture}
\end{center}
\end{tkzexample}



%====================================================================
\subsection{宏 \tkzMacro{tkz-elements}{tkzDrawCoordinates}}
该宏用 TikZ 的 \code{plot coordinates} 语法绘制一条曲线，坐标由 Lua 生成。

\medskip
\texttt{概念视角：}

Lua 的 \code{path} 是一个计算对象：它存放 Lua 阶段产生的几何数据（坐标、采样、变换、拼接等）。

TikZ 命令 \code{plot coordinates} 是一条绘图指令：它只是把收到的显式坐标画出来。

\medskip
\emph{Lua 负责计算 —— TikZ 负责绘制。}

\medskip
宏 \tkzMacro{tkz-elements}{tkzDrawCoordinates} 在这两个世界之间架起桥梁：

\begin{itemize}
\item Lua 构造并操作几何对象；
\item TikZ 渲染这份显式的坐标列表；
\item 绘图时不再做额外的几何解释。
\end{itemize}

这种严格分离保证了计算出的几何与最终图形输出之间的一致性。


\texttt{语法： } |\tkzDrawCoordinates[<options>](lua-path)|

\texttt{用法示例： }


\begin{tkzexample}[latex=.5\textwidth]
\directlua{
  z.A = point(0, 0)
  z.B = point(6, 0)
  z.C = point(1, 4)
  T.ABC = triangle(z.A, z.B, z.C)
  z.M = T.ABC.bc.mid
  z.N = T.ABC.ca.mid
  PA.path = T.ABC:path(z.M,z.N,4)
}
\begin{tikzpicture}[gridded]
\tkzGetNodes
\tkzDrawPolygon(A,B,C)
\tkzDrawCoordinates[orange,thick](PA.path)
\tkzDrawPoints(A,B,C,M,N)
\tkzLabelPoints(A,B)
\tkzLabelPoints[above](C)
\end{tikzpicture}
\end{tkzexample}


\medskip
\texttt{描述： }
把一个 Lua \code{path} 对象作为 TikZ 路径绘制。本宏通常在 Lua 计算阶段之后使用。它用显式的 \code{(x,y)} 坐标，绘制由一个 Lua 路径中的点序列所定义的折线。


\medskip
\texttt{别名： } % \let\tkzDrawPath\tkzDrawCoordinates
\label{subsec:tkzDrawPath}

\begin{quote}
\tkzMacro{tkz-elements}{tkzDrawPath[<tikz-options>](lua-path)}
\end{quote}

%====================================================================

\subsection{宏 \tkzMacro{tkz-elements}{tkzDrawPointsFromPath}}
\label{sub:macro_tkzDrawPointsFromPath}

该宏用 \tkzname{tkzDrawPoint} 绘制一条路径上的所有点，而不导出点的名字。

\medskip
\texttt{语法： }
\verb|\tkzDrawPointsFromPath[<options>](lua-path)|

\medskip
\texttt{用法示例： }

\begin{tkzexample}[latex=.5\textwidth]
\directlua{
init_elements()
LP.test = list_point()
LP.test:add(point(0, 0))
LP.test:add(point(1, 0))
LP.test:add(point(1, 1))
PA.curve = LP.test:as_path()
}
\begin{tikzpicture}
\tkzDrawCoordinates[blue](PA.curve)
\tkzDrawPointsFromPath[red,size=2](PA.curve)
\end{tikzpicture}
\end{tkzexample}

%====================================================================


%====================================================================
\subsection{宏 \tkzMacro{tkz-elements}{tkzGetPointsFromPath}}
该宏由一个 Lua 路径定义若干 TikZ 点。每个点以一个基名加索引命名。

\texttt{语法： }\verb|\tkzGetPointsFromPath[<options>](<path>,<basename>)|

\medskip
\texttt{描述： }
把一个 Lua 路径的点导出为具名的 \TeX/TikZ 点。
命名方式为 \code{<prefix>1}、\code{<prefix>2}、……

\medskip
\texttt{示例： }

\begin{tkzexample}[latex=.5\textwidth]
 \directlua{
  init_elements()
  PA.g = path({ "(0, 0)", "(1, 0)","(1, 1)" })
}
\begin{tikzpicture}
  \tkzGetPointsFromPath(PA.g,A)
  \tkzDrawPolygon(A_1,A_2,A_3)
  \tkzDrawPoints(A_1,A_2,A_3)
\end{tikzpicture}
\end{tkzexample}

%====================================================================



\subsection{宏 \tkzMacro{tkz-elements}{tkzGetPointFromPath}}
该宏取出一个 Lua 路径的第 i 个顶点，并把它定义为一个给定名字的 TikZ 点。

\texttt{语法： }\verb|\tkzGetPointFromPath(<PathLua>,<index>){<PointName>}|

\texttt{参数： }

\verb|<PathLua>|：一个 Lua \verb|path| 对象的标识符。
\verb|<index>|：要取回的顶点索引（TeX 数或数值宏）。
\verb|<PointName>|：要创建的 TikZ 点的名字。
\medskip


\texttt{用法示例：}

\begin{mybox}
\begin{verbatim}
% Get the 3rd vertex of PA.A and name it P3
\tkzGetPointFromPath(PA.A,3){P3}
% With counting and a loop
\tkzPathCount(PA.A){N}
\foreach \i in {1,...,\N}{
\expandafter\tkzGetPointFromPath\expandafter(PA.A,\i){P\i}
}
\end{verbatim}
\end{mybox}
%====================================================================




\subsection{宏 \tkzMacro{tkz-elements}{tkzDrawSegmentsFromPaths}}
\label{sub:macro_tkzDrawSegmentsFromPaths}

\texttt{描述： }
绘制编码在一条或多条 Lua 路径中的线段（点对或连续点，取决于所采用的约定）。

\texttt{语法： }\verb|\tkzDrawSegmentsFromPaths[<options>](<lua-path-start>,<lua-path-end>)|

\texttt{示例： }

\begin{tkzexample}[vbox]
\directlua{
 init_elements()
 z.A = point(0, 0)
 z.B = point(3, 0)
 C.AB = circle(z.A, z.B)
 PA.c = path()
 PA.u = path()
 PA.v = path()
 for i = 0, 18, 2 do
   local angle = i * 2 * math.pi / 18
   local p = point(polar(8, angle))
   local Lu, Lv = C.AB:tangent_from(p)
   local u , v  = Lu.pb, Lv.pb
   PA.c:add_point(p)
   PA.u:add_point(u)
   PA.v:add_point(v)
 end}

\begin{center}
 \begin{tikzpicture}[scale =.6]
  \tkzGetNodes
  \tkzDrawSegmentsFromPaths[draw,red](PA.c,PA.u)
  \tkzDrawSegmentsFromPaths[draw,red](PA.c,PA.v)
  \tkzDrawCircle[blue](A,B)
\end{tikzpicture}
 \end{center}
\end{tkzexample}
%====================================================================
\subsection{宏 \tkzMacro{tkz-elements}{tkzDrawCirclesFromPaths}}
\label{sub:macro_tkzDrawCirclesFromPaths}
该宏绘制由两条路径定义的圆：一条给出圆心，另一条给出每个圆所经过的点。

\medskip
\texttt{语法： }\verb|\tkzDrawCirclesFromPaths[<options>](<PA.center>,<PA.through>)|

\medskip
\texttt{用法示例： }

\begin{mybox}
\begin{verbatim}
\tkzDrawCirclesFromPaths[blue](PA.O, PA.A)
\end{verbatim}
\end{mybox}

更多示例见在 \href{http://altermundus.fr}{altermundus.fr} 上发布的文档《欧几里得几何》。

\medskip
\begin{tkzexample}[vbox]
\directlua{
 init_elements()
 z.A = point(-4, -4)
 z.B = point(4, -4)
 L.AB = line(z.A, z.B)
 C.AB = circle(z.A, z.B)
 S.AB = L.AB:square()
 _, _, z.C, z.D = S.AB:get()
 z.O = S.AB.ac.mid
 z.a = S.AB.ab.mid
 z.c = S.AB.cd.mid
 z.o1 = tkz.midpoint(z.a, z.O)
 z.o = tkz.midpoint(z.c, z.O)
 z.b = S.AB.bc.mid
 z.g = (z.o1 - z.O) + (z.b - z.O)
 z.o2 = intersection(line(z.c, z.g), line(z.O,z.b))
 z.o3 = (z.o - z.O) + (z.o2 - z.O)
 z.j = intersection(line(z.c, z.b), line(z.o2,z.o3))
 local C2 = circle(z.O, z.a)
 local C1 = circle(z.o, z.c)
 local C3 = circle(z.o3, z.j)
 PA.pc, PA.pt, n = C1:CCC(C2, C3)
}

\begin{center}
\begin{tikzpicture}[scale=1]
\tkzGetNodes
\tkzDrawPolygon[cyan](A,B,C,D)
\tkzDrawCircles(O,a o,c o1,a o2,b o3,j)
\tkzDrawCirclesFromPaths[draw,red](PA.pc,PA.pt)
\end{tikzpicture}
\end{center}
\end{tkzexample}


%===============================================================

\subsection{宏 \tkzMacro{tkz-elements}{tkzDrawFromPointToPath}}
该宏从一个给定的 TikZ 点出发，向一个 Lua \emph{path} 对象的每个点画一组线段。

\medskip
\texttt{语法： }\verb|\tkzDrawFromPointToPath[<options>](<point-tikz>,<lua-path>)|

\medskip
\texttt{参数： }

\verb|<options>|（可选）：pgf/tikz 绘图选项（颜色、线宽等）。

\verb|<point-tikz>|：一个已定义的 TikZ 点的名字。

\verb|<lua-path>|：一个 Lua \verb|path| 对象的标识符。

\medskip
\texttt{用法示例： }

\medskip
\begin{tkzexample}[vbox]
\directlua{
 init_elements()
 z.O = point(0, 0)
 z.A = point(5, 1)
 z.Ta = point(8, 2)
 C.A = circle(z.A, z.Ta)
 L.OA = line(z.O,z.A)
 z.x, z.y = intersection(C.A, L.OA,{near=z.O})
 C.A.through = z.y
 z.m = tkz.midpoint(z.O, z.x)
 z.n = tkz.midpoint(z.O, z.y)
 z.w = tkz.midpoint(z.m, z.n)
 PA.A = path()
 PA.c = path()
 for t = 0, 1 - 1e-12, 0.1 do
   PA.A:add_point(C.A:point(t))
   local m = tkz.midpoint(z.O, C.A:point(t))
   PA.c:add_point(m)
 end
}
 \begin{center}
  \begin{tikzpicture}[scale=.8]
   \tkzGetNodes
   \tkzDrawCircle(A,Ta)
   \tkzDrawLine(O,A)
   \tkzDrawPoints(O,A,m,n,w,x,y)
   \tkzLabelPoints(O,A,m,n,w)
   \tkzDrawCircle[red](w,n)
   \tkzDrawPointsFromPath[red](PA.c)
   \tkzDrawPointsFromPath[blue](PA.A)
   \tkzDrawSegmentsFromPaths[draw,teal](PA.c,PA.A)
   \tkzDrawFromPointToPath[orange](O,PA.A)
 \end{tikzpicture}
  \end{center}
\end{tkzexample}



\subsection{阿基米德螺线：一个完整示例}

本示例借助本节介绍的若干宏构建而成。它的灵感来自 Jean-Marc Desbonnez 用 tkz-euclide 创建的一个示例。

\vspace{1em}
\begin{tkzexample}[vbox]
\directlua{
  z.O = point(0, 0)
  z.A = point(5, 0)
  L.OA = line(z.O, z.A)
  PA.spiral = path()
  PA.center = path()
  PA.through = path()
  PA.ray = path()
  for i = 1, 48 do
      local k = i / 48
      z["P"..i] = L.OA:point(k)
      local angle = i * tkz.tau / 48
      z["r"..i] = point(polar(5, angle))
      ray = line(z.O, z["r"..i])
      circ = circle(z.O, z["P"..i])
      z["X"..i], _ = intersection(ray, circ, { near = z["r"..i] })
      PA.spiral:add_point(z["X"..i])
      PA.center:add_point(z.O)
      PA.through:add_point(z["P"..i])
      PA.ray:add_point(z["r"..i])
  end}
\begin{center}
  \begin{tikzpicture}[scale = .8]
  \tkzGetNodes
  \tkzDrawCircle(O,A)
  \tkzDrawCirclesFromPaths[draw,
      orange!50](PA.center,PA.through)
  \tkzDrawSegmentsFromPaths[draw,teal](PA.center,PA.ray)
  \tkzDrawPath[red,thick,smooth](PA.spiral)
  \tkzDrawPointsFromPath[blue, size=2](PA.spiral)
  \end{tikzpicture}
\end{center}

\end{tkzexample}



\endinput